% RESEARCH_DRAFT_SHA256: 71c1ef59b27356baa1d1ef4aad65ead7d13e9617ebd5704ef849962d2706beaa
% Layout and source-reference wrappers expanded for the web renderer.
\documentclass[11pt]{article}
\usepackage[margin=25mm]{geometry}
\usepackage{amsthm}
\usepackage{fontspec,unicode-math}
\setmainfont{Latin Modern Roman}
\setmathfont{Latin Modern Math}
\usepackage{amsmath,amssymb,mathtools,xurl,hyperref}
\hypersetup{colorlinks=true,urlcolor=blue,linkcolor=blue}
\setcounter{secnumdepth}{0}
\setlength{\parindent}{0pt}
\setlength{\parskip}{5pt}
\setlength{\emergencystretch}{3em}
\newcommand{\sref}[2]{\hyperlink{source:#1:#2}{[S#2]}}
\newcommand{\cataloguescope}[1]{\paragraph{Source record.}#1}
\newcommand{\floor}[1]{\left\lfloor#1\right\rfloor}
\newtheorem{proposition}{Proposition}
\newtheorem{lemma}[proposition]{Lemma}
\begin{document}
\paragraph{Research manuscript.} OpenAI \texttt{gpt-6-astra}, reasoning \texttt{ultra}; 11 October 2026.
\paragraph{Preserved source body.} The complete catalogue mathematical text and sources follow with layout and source links expanded, including source qualifications and its historical preparation notes. New research begins at the section headed Research attempt.

% UnsolvedMath ID 30000237; source catalogue code OWR-831-009
\section{Hardness of Decoding Random Linear Codes}\label{30000237}
{\small UnsolvedMath ID 30000237 · Computer Science}
\subsection{Definitions and mathematical statement}

\paragraph{Shared notation.}Standard number systems are $\mathbb N,\mathbb Z,\mathbb Q,\mathbb R,\mathbb C$. Specialized notation and index ranges follow the cited source and the definitions below.


A random linear code may be specified by a randomly chosen generator or parity-check matrix over a finite field. Decoding seeks a codeword close to a received word, and average-case hardness depends on this sampling distribution.
\paragraph{Record scope and qualifications.}
The catalogue title is retained as metadata. Named certificates, contracts, bounds or frameworks in that title are not asserted as results of the original source extract.
The alphabet, rate, noise distribution, decoding radius, and search/decision model are not specified in the supplied question.

Source-record limitation: The code distribution, decoding radius/noise regime, and success target are unspecified.
\paragraph{Statement recorded in the supplied source.}
Determine the computational hardness of decoding from random linear codes. \href{https://doi.org/10.4171/owr/2005/26}{[S2]}
\paragraph{Source status.}
The supplied catalogue labels this source record partially solved; that metadata was checked against the supplied snapshot on 2026-10-11. The label is preserved for provenance and is not an independent verification of present open status or mathematical correctness. The catalogue literature assessment is dated 2026-08-21.
\subsection{Short English statement}
Classify the computational difficulty of decoding randomly generated linear codes.
\subsection{Sources}
\begin{itemize}
\item[S1] Ulam AI and the UnsolvedMath Contributors, UnsolvedMath: A Curated Collection of Open Mathematics Problems, record 30000237 (OWR-831-009). Catalogue curation and metadata: CC BY 4.0. The source extract is retained; definitions, English summaries and qualifications are editorial additions.
\url{https://huggingface.co/datasets/ulamai/UnsolvedMath}.
\item[S2] Complexity Theory (2006), Oberwolfach Reports, DOI 10.4171/owr/2005/26, pp. 1431--1504
\url{https://doi.org/10.4171/owr/2005/26}.
\item[S3] Random-code decoding results
\url{https://eccc.weizmann.ac.il/report/2010/003/}.
\item Source-specific Oberwolfach Reports licences, including CC BY-SA 4.0 where applicable, are retained; see THIRD\_PARTY\_NOTICES.md.
\end{itemize}
\paragraph{Preparation and review.}
The original source and UnsolvedMath contributors retain attribution for the source extract. Codex (GPT-6; exact serving revision not exposed) assisted with definitions, English summaries, scope notes and formatting. No new human mathematical authorship or independently verified solution is asserted. Permanent source ID 30000237 and code OWR-831-009 checked against the complete supplied registry and current repository. Distinct upstream records, including repeated questions and broad or damaged extracts, are retained for complete coverage of the source CS category. A repeated record is not a new independent conjecture. Complete mathematical-equivalence and current-literature review has not been performed. Preparation and final review follow contribution prompts C and D, with source gaps disclosed. Structural and rendering checks do not establish a new theorem or open status.
\label{end:30000237}

\section{Research attempt}
\subsection{Literature, source context and scope}
The full original catalogue body and source qualifications are preserved above. The relevant primary workshop passages were read in the cached publisher report. Historical claims and subsequent attributed results are distinguished from the mathematical work below. A restricted example, standard lemma or known theorem is not counted as new progress on the entire original target. Missing assumptions and repeated or rhetorical extracts are not silently turned into different solved problems.
Codex (GPT-6; exact serving revision not exposed) supplied the auxiliary proof route used here. OpenAI \texttt{gpt-6-astra}, reasoning \texttt{ultra}, independently worked through its mathematical steps, wrote this exposition, and performed the finite checks. This is model collaboration, not independent human verification.

\subsection{A random-code identifiability calculation}
To separate sample information from computational hardness, specify one model: $Q\in\mathbb F_q^{m\times n}$ has independent uniform entries, $s\in\mathbb F_q^n$ is fixed, and $b=Qs+e$, where the errors are independent of $Q$, independent across coordinates, and have common zero probability $\theta>1/q$. Their distribution on nonzero symbols is otherwise arbitrary. We examine exhaustive nearest-codeword recovery of $s$.
\begin{lemma}
If independent variables $X_i\in[0,1]$ have sum mean $\mu$, then for $a>0$, either tail at distance $ma$ from $\mu$ has probability at most $e^{-2ma^2}$.
\end{lemma}
\begin{proof}
For one such variable let $g(t)=\log\mathbb E e^{tX}$. The second derivative is the variance under the tilted probability law. Since $\mathbb E(X-1/2)^2\le1/4$ under every law supported on $[0,1]$, this variance is at most $1/4$. With $g(0)=0$ and $g'(0)=\mathbb EX$, twice integrating gives $\log\mathbb E e^{t(X-\mathbb EX)}\le t^2/8$ for every real $t$. Independence and exponential Markov give an upper-tail bound $\exp(-tma+mt^2/8)$ for $t>0$. Choosing $t=4a$ proves the result; applying the same argument to $-\sum(X_i-\mathbb EX_i)$ proves the lower tail.
\end{proof}
\begin{proposition}
Except with probability at most $q^n\exp[-m(\theta-1/q)^2/2]$, the true $s$ is the unique message minimizing the Hamming discrepancy between $b$ and $Qs'$.
\end{proposition}
\begin{proof}
The true discrepancy is a sum of independent Bernoulli variables of mean $1-\theta$. For each fixed $s'\ne s$, the nonzero linear form given by each row of $Q(s'-s)$ is uniform over $\mathbb F_q$, independently across rows and independently of the errors. Thus the wrong discrepancy is a sum of independent Bernoulli variables of mean $1-1/q$. Use the intermediate threshold
\[
\tau=m\bigl[1-(\theta+1/q)/2\bigr].
\]
The probability that the true discrepancy is at least $\tau$, or that a particular wrong discrepancy is at most $\tau$, is bounded by $\exp[-m(\theta-1/q)^2/2]$ by the lemma. A union bound over one true and $q^n-1$ wrong events proves the assertion. The strict inequalities outside these events also exclude ties, whether or not $\tau$ is integral.
\end{proof}
\subsection{Verification and remaining gaps}
A seeded instance with $q=3,n=2,m=100,\theta=3/4$ was checked against all nine messages; the theoretical union bound was evaluated as well. This is a polynomial sample bound when the parameters permit a fixed positive gap $\theta-1/q$, but the described search enumerates $q^n$ candidates. The attempted inference from statistical distinguishability to efficient decoding stops at that enumeration. Conversely, worst-case hardness of a decoding problem does not by itself imply average-case hardness under this independent random-code/noise distribution. No such reduction, efficient decoder or general hardness classification is proved. This is an auxiliary known-style calculation for a specified ensemble, not completion of the report's broad decoding-complexity question.

\section{References and source attribution}
Primary workshop report: \url{https://doi.org/10.4171/owr/2005/26}. Relevant passages read: 2005/26 p. 1453, learning modulo a prime and decoding random linear codes.

All retrievals were checked on 11 October 2026. Original question attribution remains in the preserved Sources section. External theorems are attributed and not claimed as new discoveries or fully verified merely by citation.


\paragraph{Integration review.} Codex (GPT-6; exact serving revision not exposed) reviewed the written arguments and expanded layout and source-reference wrappers. Repository metadata and the production web renderer were checked using Python 3.13 and JavaScript with MathJax 3.2.2. This model-assisted review is not independent human proof verification; no PDF compilation or formal proof checking is asserted.

\section{Completion Estimate}
\noindent\textbf{COMPLETION: 0\%}
No substantiated new progress toward the entire original target is claimed. The worked results and failed route are delimited above; a known result, illustrative model, or auxiliary obstruction is not counted as original-target completion.

\section{Final status and attribution}
\textbf{Attempt status: unresolved.} This describes this research attempt, not a new ruling on the source catalogue's status or on external literature claims.
The mathematics and this self-audit were generated by OpenAI \texttt{gpt-6-astra}, reasoning \texttt{ultra}, on 11 October 2026. Codex (GPT-6; exact serving revision not exposed) orchestrated the batch and contributed formatting/repository review. No human mathematical authorship, independent human proof checking, or formal verification is claimed. The historical preparation paragraphs in the reproduced catalogue body concern that body only.

Contribution prompt A was applied to the complete definition and supplied source record, followed by the adversarial checks of prompt B. Local reproduction files are \texttt{batch016/checks.py}, \texttt{computations.json}, and the per-ID audit JSON. They were not public links at preparation time. Numerical checks are finite and do not prove unrestricted results. No LaTeX engine was found locally; compilation was not performed by this research agent.

\end{document}
